% year: תש"פ
% type: מבחן
% semester: ב
% moed: ב
% number: 1ג
% points: 5
% categories: func-analysis, func-limits
% source: exams/מבחנים/תשפ/מועד ב סמסטר ב תשפ.pdf

\begin{question}
(5 נק') לאילו ערכים של הפרמטר $a$ לגרף של פונקציה $f(x)=\frac{x^2-9}{x+a}$ אין אסימפטוטה אנכית?
\end{question}

\begin{hint}
מתי השורש של המכנה הוא גם שורש של המונה?
\end{hint}

\begin{solution}
הפונקציה רציפה בכל נקודה שבה $x\neq -a$, ולכן אסימפטוטה אנכית יכולה להיות רק ב-$x=-a$. ב-$x\to-a$ המכנה שואף ל-$0$, ולכן:
\begin{itemize}
  \item אם המונה אינו מתאפס ב-$-a$, כלומר $a^2-9\neq0$, אז $|f(x)|\to\infty$ כאשר $x\to-a$, ו-$x=-a$ היא אסימפטוטה אנכית.
  \item אם $a^2=9$, כלומר $a=3$ או $a=-3$: עבור $a=3$, $f(x)=\frac{(x-3)(x+3)}{x+3}=x-3$ לכל $x\neq-3$, ולכן $\lim_{x\to-3}f(x)=-6$ סופי (נקודת אי-רציפות סליקה), ואין אסימפטוטה אנכית. באופן דומה עבור $a=-3$, $f(x)=x+3$ לכל $x\neq3$ ו-$\lim_{x\to3}f(x)=6$.
\end{itemize}
לכן אין אסימפטוטה אנכית בדיוק כאשר $\boxed{a=3 \text{ או } a=-3}$.
\end{solution}
